\documentclass[12pt,a4paper]{article}
\usepackage[utf8]{inputenc}
\usepackage[vietnamese]{babel}
\usepackage[T5]{fontenc}
\usepackage{amsmath,amssymb,amsfonts}
\usepackage{fancyhdr}
\usepackage{geometry}
\usepackage{tcolorbox}

\geometry{
	a4paper,
	left=20mm,
	right=20mm,
	top=20mm,
	bottom=20mm
}

\pagestyle{fancy}
\fancyhf{}
\lhead{\small\textit{Tài liệu Hình học Chuyên đề -- Bổ đề cốt lõi}}
\rhead{\small\textbf{Integra Typesetter}}
\rfoot{\small Trang \thepage}
\renewcommand{\headrulewidth}{0.4pt}
\renewcommand{\footrulewidth}{0.4pt}

\begin{document}

\begin{center}
	{\Large \textbf{CÁC BỔ ĐỀ CỐT LÕI HÌNH HỌC PHẲNG}}\\[2mm]
	{\small \textit{(Ứng dụng trong các bài toán Hình học Vận dụng cao / Chuyên Toán)}}
\end{center}
\vspace{4mm}

\begin{tcolorbox}[colback=blue!5!white,colframe=blue!75!black,title=\textbf{Bổ đề 1: Hệ thức khoảng cách Trực tâm và Cặp tia đẳng giác}]
Cho tam giác $ABC$ nhọn nội tiếp đường tròn $(O; R)$ có trực tâm $H$. Khi đó:
\begin{enumerate}
    \item $\widehat{BAH} = \widehat{CAO} = 90^\circ - \widehat{B}$ (Tia $AH$ và $AO$ đẳng giác trong góc $\widehat{BAC}$).
    \item $AH = 2R \cos \widehat{BAC}$. Đặc biệt, nếu $\widehat{BAC} = 60^\circ$ thì $AH = R = AO$.
\end{enumerate}
\end{tcolorbox}

\vspace{2mm}

\begin{tcolorbox}[colback=teal!5!white,colframe=teal!75!black,title=\textbf{Bổ đề 2: Đường đối cực của một điểm đối với đường tròn qua Tứ giác toàn phần}]
Cho tứ giác toàn phần $AXDY.PH'$ nội tiếp đường tròn $(O)$ (với $AD$ là đường kính, $X = PA \cap (O)$, $Y = PD \cap (O)$). Khi đó, giao điểm của các cặp cạnh đối và giao điểm của hai đường chéo tạo thành tam giác tự cực. Đặc biệt, đường đối cực của điểm $P$ đối với $(O)$ đi qua giao điểm của $XY$ và $AD$, đồng thời đi qua trực tâm tương ứng $H'$.
\end{tcolorbox}

\vspace{2mm}

\begin{tcolorbox}[colback=orange!5!white,colframe=orange!75!black,title=\textbf{Bổ đề 3: Định lý La Hire}]
Trong mặt phẳng cho đường tròn $(O)$. Điểm $M$ nằm trên đường đối cực của điểm $N$ khi và chỉ khi điểm $N$ nằm trên đường đối cực của điểm $M$.
\end{tcolorbox}

\end{document}
